\section*{अध्याय \ref{ch:g12:mechanical-waves} --- यांत्रिक तरंगें}

\begin{solution}{exo:g12:mechanical-waves:1}
(क) $v = 320/40 = \qty{8.0}{m/s}$। (ख) कोई दर्शक घेरे के साथ नहीं चलता: चलता विक्षोभ (खड़े होने का
कर्म) है, भीड़ नहीं। (ग) $12/1.5 = \qty{8.0}{m/s}$ --- संगत।
\end{solution}

\begin{solution}{exo:g12:mechanical-waves:2}
(क) \omterm{def:g12:mechanical-waves:transverse}{अनुप्रस्थ}; (ख) \omterm{def:g12:mechanical-waves:transverse}{अनुदैर्घ्य}; (ग) \omterm{def:g12:mechanical-waves:transverse}{अनुदैर्घ्य};
(घ) \omterm{def:g12:mechanical-waves:transverse}{अनुप्रस्थ}; (ङ) P \omterm{def:g12:mechanical-waves:transverse}{अनुदैर्घ्य}, S \omterm{def:g12:mechanical-waves:transverse}{अनुप्रस्थ}।
\end{solution}

\begin{solution}{exo:g12:mechanical-waves:3}
(क) $d = 340 \times 4.0 \approx \qty{1.4}{km}$। (ख) $1500 \times 4.0 = \qty{6.0}{km}$। (ग) जिस माध्यम को पार किया गया है उसमें
\omterm{def:g12:mechanical-waves:speed}{संचरण चाल}।
\end{solution}

\begin{solution}{exo:g12:mechanical-waves:4}
$\lambda = v/f$: हवा में $340/440 = \qty{0.77}{m}$, पानी में $1500/440 = \qty{3.4}{m}$। आवृत्ति वही रहती है (उसे
स्रोत तय करता है); $v$ और $\lambda$, दोनों बदल जाते हैं।
\end{solution}

\begin{solution}{exo:g12:mechanical-waves:5}
$\lambda = \qty{12}{cm}$; $T = \lambda/v = 0.12/0.30 = \qty{0.40}{s}$; $f = 1/T = \qty{2.5}{Hz}$।
\end{solution}

\begin{solution}{exo:g12:mechanical-waves:6}
(क) $v = 2.4/0.30 = \qty{8.0}{m/s}$। (ख) पहले चित्र से शिखर को अभी $\qty{6.0}{m}$ और चलना है: यानी उसके
$t = 6.0/8.0 = \qty{0.75}{s}$ बाद। (ग) रिबन विशुद्ध \omterm{def:g12:mechanical-waves:transverse}{अनुप्रस्थ} दिशा में उठता है और फिर गिर जाता है,
और $x = \qty{5.0}{m}$ पर ही बना रहता है: रस्सी चलती नहीं।
\end{solution}

\begin{solution}{exo:g12:mechanical-waves:7}
(क) पास आने की चाल \qty{4.0}{m/s}, दूरी \qty{4.0}{m}: $t = \qty{1.0}{s}$। (ख) अध्यारोपण: $3.0 + 2.0 = \qty{5.0}{cm}$।
(ग) $3.0 - 2.0 = \qty{1.0}{cm}$। (घ) दोनों स्पंदें साबुत निकलकर ज्यों की त्यों चलती रहती हैं,
मानो वे कभी मिली ही न हों।
\end{solution}

\begin{solution}{exo:g12:mechanical-waves:8}
(क) पतली वाली: बराबर \omterm{def:g10:inertia:inventory}{तनाव} पर प्रति \omterm{def:g10:orders-of-magnitude:unit}{मीटर} कम द्रव्यमान
का अर्थ है तेज़ तरंग। (ख) वह बढ़ जाती है (\omterm{def:g10:inertia:inventory}{तनाव} अधिक)।
(ग) नहीं: चाल केवल माध्यम पर निर्भर करती है, स्पंद के नाप पर नहीं।
\end{solution}

\begin{solution}{exo:g12:mechanical-waves:9}
पटरी: $1800/5900 \approx \qty{0.31}{s}$; हवा: $1800/340 \approx \qty{5.3}{s}$; अंतर $\approx \qty{5.0}{s}$।
\end{solution}

\begin{solution}{exo:g12:mechanical-waves:10}
$\lambda = \qty{24}{cm}$ के साथ: (क) $2\lambda$, \omterm{def:g12:mechanical-waves:phase}{समकला}; (ख) $\lambda/2$, \omterm{def:g12:mechanical-waves:phase}{विपरीत कला};
(ग) $\tfrac32\lambda$, \omterm{def:g12:mechanical-waves:phase}{विपरीत कला}; (घ) $1.25\lambda$, कोई नहीं; (ङ) $\lambda/4$, कोई नहीं।
\end{solution}

\begin{solution}{exo:g12:mechanical-waves:11}
$\lambda_P = 6.0 \times 2.0 = \qty{12}{km}$; $\lambda_S = 3.5 \times 2.0 = \qty{7.0}{km}$ --- यानी किसी इमारत के नाप का सैकड़ों गुना। $\lambda$ से कहीं
पास बैठे बिंदु लगभग \omterm{def:g12:mechanical-waves:phase}{समकला} में रहते हैं, इसलिए पूरा मोहल्ला एक साथ
हिलता है।
\end{solution}

\begin{solution}{exo:g12:mechanical-waves:12}
(क) $T = \qty{0.50}{s}$, $f = \qty{2.0}{Hz}$। (ख) $\lambda = \qty{0.75}{m}$। (ग) $v = \lambda/T = 0.75/0.50 = \qty{1.5}{m/s}$। (घ) $t = T = \qty{0.50}{s}$ पर। (ङ) $1.875/0.75 = 2.5$
\omterm{def:g12:mechanical-waves:wavelength}{तरंगदैर्घ्य} $= (2 + \tfrac12)\lambda$: \omterm{def:g12:mechanical-waves:phase}{विपरीत कला}।
\end{solution}

\begin{solution}{exo:g12:mechanical-waves:13}
(क) $T = 80/10 = \qty{8.0}{s}$, $f = \qty{0.125}{Hz}$। (ख) $v = \lambda/T = 120/8.0 = \qty{15}{m/s} = \qty{54}{km/h}$। (ग) $t = 900/15 = \qty{60}{s}$। (घ) नहीं: तरंगें पदार्थ
नहीं ढोतीं --- लंगर डाला प्लव केवल ऊपर-नीचे होता रहता है।
\end{solution}

\begin{solution}{exo:g12:mechanical-waves:14}
(क) \qty{6.0}{m} की दूरी पर \qty{6.0}{m/s} से पास आते हुए: $t = \qty{1.0}{s}$; विपरीत विस्थापन कट जाते
हैं --- रस्सी पल भर के लिए चपटी हो जाती है। (ख) नहीं: उसके बिंदु चल रहे हैं
(\omterm{def:g12:mechanical-waves:transverse}{अनुप्रस्थ} वेग जुड़ते हैं, कटते नहीं); पूरी ऊर्जा गतिज है।
(ग) दोनों स्पंदें फिर प्रकट होकर ज्यों की त्यों चलती रहती हैं। (घ) चपटी रस्सी
वेग ढो रही है, यानी ऊर्जा: तरंग पल भर के लिए आकृति में नहीं,
गति में संचित है।
\end{solution}

\begin{solution}{exo:g12:mechanical-waves:15}
(क) $2\pi r$: \qty{3.1}{m} और \qty{50}{m} --- $\times 16$। (ख) प्रति \omterm{def:g10:orders-of-magnitude:unit}{मीटर}
ऊर्जा $16$ से गिरती है, इसलिए आयाम भी घटता है
(ऊर्जा आयाम के साथ बढ़ती है)। (ग) सीधे अग्र लंबे नहीं होते:
इसलिए केवल अवशोषण बचता है। (घ) वहाँ माध्यम स्वयं सक्रिय है: हर दर्शक अपनी
माँसपेशियों की ऊर्जा से खड़ा होता है और तरंग को फिर से भर
देता है।
\end{solution}

\begin{solution}{pb:g12:mechanical-waves:1}
\textbf{1.} P सफ़र की दिशा में दबाती हैं: \omterm{def:g12:mechanical-waves:transverse}{अनुदैर्घ्य}; S बग़ल में
हिलाती हैं: \omterm{def:g12:mechanical-waves:transverse}{अनुप्रस्थ}।

\textbf{2.} तेज़ होने के कारण P हर जगह पहले पहुँचती है। $t_P = 120/6.0 = \qty{20}{s}$; $t_S = 120/3.5 \approx \qty{34}{s}$।

\textbf{3.} $\Delta t = d/v_S - d/v_P = d\,(1/v_S - 1/v_P)$; और \qty{120}{km} के लिए: $34.3 - 20 \approx \qty{14}{s}$।

\textbf{4.} $d$ के लिए हल कीजिए: अंतराल के प्रति \omterm{def:g10:orders-of-magnitude:unit}{सेकंड} $d = \dfrac{v_P v_S}{v_P - v_S}\,\Delta t
= \dfrac{6.0 \times 3.5}{2.5}\,\Delta t = \qty{8.4}{km}$।

\textbf{5.} $\Delta t \propto d$: यानी किसी एक स्टेशन का अपना अभिलेख ही, आरंभ-समय जाने
बिना, अधिकेंद्र से उसकी दूरी दे देता है।

\textbf{6.} $d_A = 8.4 \times 25 = \qty{2.1e2}{km}$।

\textbf{7.} अंतराल दूरी देता है, दिशा नहीं: अधिकेंद्र A के चारों ओर त्रिज्या
$d_A$ वाले वृत्त पर कहीं भी हो सकता है।

\textbf{8.} $d_B = 8.4 \times 15 = \qty{1.3e2}{km}$; $d_C = 8.4 \times 32 \approx \qty{2.7e2}{km}$।

\textbf{9.} दो वृत्त दो बिंदुओं पर मिलते हैं; तीसरा उनमें से एक चुन लेता है:
तीनों का साझा प्रतिच्छेद ही अधिकेंद्र है।

\textbf{10.} $t_P = 210/6.0 = \qty{35}{s}$: भूकंप $09{:}14{:}32 - \qty{35}{s} = 09{:}13{:}57$ पर आरंभ हुआ।

\textbf{11.} $\qty{200}{m/s} = \qty{720}{km/h}$ --- यानी यात्री जेट विमान की रफ़्तार।

\textbf{12.} $t = \num{6.0e5}/200 = \qty{3.0e3}{s} = \qty{50}{min}$।

\textbf{13.} $t_P = 600/6.0 = \qty{100}{s}$। चेतावनी: $3000 - 100 = \qty{2.9e3}{s} \approx \qty{48}{min}$।

\textbf{14.} $\lambda = vT = 200 \times 900 = \qty{1.8e5}{m} =
\qty{180}{km}$: \qty{180}{km} पर फैली मुश्किल से एक \omterm{def:g10:orders-of-magnitude:unit}{मीटर} की चढ़ाई
--- जहाज़ कई मिनटों में एक अगोचर ढाल चढ़ जाता है।

\textbf{15.} $\lambda = 10 \times 900 = \qty{9.0}{km}$: तरंग बीस गुना सिकुड़ जाती है, उसकी
ऊर्जा कहीं छोटी लंबाई में ठुँस जाती है, इसलिए उसकी ऊँचाई बढ़नी ही
है --- और सूनामी तट पर आकर तन जाती है।

\textbf{16.} $2\pi r$: \qty{63}{km}, फिर $\approx \qty{1.3e3}{km}$ --- $\times 21$।

\textbf{17.} वही ऊर्जा $21\times$ गुना लंबे अग्र पर बँट जाती है, इसलिए
हानिरहित भूपर्पटी में भी आयाम गिरता है; और असली भूपर्पटी उसमें
अवशोषण भी जोड़ देती है।

\textbf{18.} S तरंगें, यानी \omterm{def:g12:mechanical-waves:transverse}{अनुप्रस्थ} अपरूपण, द्रव को पार
नहीं कर सकतीं: S-तरंग की यह छाया दिखाती है कि (बाहरी) \omterm{def:g10:refraction:fiber}{क्रोड} द्रव है।

\textbf{19.} $\num{1e15}/\num{1.5e9} \approx \num{6.7e5}$: यानी पूरी रफ़्तार पर दौड़ती लगभग सात लाख रेलगाड़ियाँ।

\textbf{20.} भूकंप $09{:}13{:}57$ पर आया, समुद्र में तीनों वृत्तों के कटान-बिंदु
पर, स्टेशन A से \qty{210}{km} दूर। बंदरगाह का P-तरंग वाला चेतावनी-यंत्र
आरंभ के \qty{100}{s} बाद बजा, यानी सूनामी के पहुँचने से लगभग \qty{48}{min} पहले।
\end{solution}
